\scanpage{36}
\begin{center}
UNIVERSITÉ DE PARIS

\figure{\begin{tikzpicture}\draw (-.4,0)--(-.06,0);\fill (0,0) circle (.7pt);\draw (.06,0)--(.4,0);\end{tikzpicture}}

ÉCOLE NORMALE SUPÉRIEURE

\figure{\begin{tikzpicture}\draw (-.4,0)--(-.06,0);\fill (0,0) circle (.7pt);\draw (.06,0)--(.4,0);\end{tikzpicture}}

45, RUE D'ULM

ODÉON 06$\cdot$45

\figure{\begin{tikzpicture}\draw (-.4,0)--(-.06,0);\fill (0,0) circle (.7pt);\draw (.06,0)--(.4,0);\end{tikzpicture}}
\end{center}

\deleted{Soit $E$ \deleted{un ensemble simplicial} un C-W-Complexe, $G$ un groupe discret, opérant sur $E$ par des transformations qui préservent la décomposition cellulaire, $B=E/G$ l'\deleted{ensemble} C-W-Complexe quotient de $E$}

\uline{Lemme 1} : $\forall x$ sommet de $E$, $\beta$ chemin de $B$ tel que $\unsure{d_1|\partial_1}(\beta)=p(x)$,
\[
\exists\gamma\text{ chemin de }E\text{ tel que }\partial_1(\gamma)=x\text{ et }p(\gamma)=\beta
\]
\uuline{Dém} Il suffit de le voir si $\beta$ est une arête. Soit $\lambda$ une arête de $B$ telle que $p\lambda=\beta$. On a \deleted{$\unsure{\partial}$} $p(\partial_1\lambda)=\partial_1(p\lambda)=\partial_1\beta=p(x)$, d'où $\exists g\deleted{\unsure{g_0}}\in G$, $x=g.\partial_1\lambda$, et $\gamma=g\lambda$ répond à la question

\uline{Lemme 2} : l'image de $p_*:\pi_1(E)\to\pi_1(B)$ est un sous groupe distingué de $\pi_1(B)$

\uuline{Dém} : Soit $\gamma$ un lacet de \deleted{$\unsure{B}$} $E$, $\beta$ un lacet de $B$, $\lambda$ un chemin de $E$ tel que $\partial_1(\lambda)=x_0$ et $p_*(\lambda)=\beta$. On a $p(\partial_0\lambda)=\partial_0\beta=b_0=px_0$, d'où $\partial_0\lambda=gx_0$, et le lacet $(\lambda,g\gamma,-\lambda)$ se projette en $(\beta,p(\gamma),-\beta)$ ce qui démontre le lemme 2.

\uline{Définition de $\chi:G\to\pi_1(B)/\operatorname{Im}\pi_1(E)$ quand $E$ connexe}.

Pour Tout $g\in G$, soit $\gamma$ un chemin de $E$ tel que $\partial_1(\gamma)=x_0$, $\partial_0(\gamma)=gx_0$. \deleted{\unsure{$\chi(g)$ la classe de}} $p(\gamma)$ est un lacet de $B$, soit $\chi(g)$ sa classe dans $\pi_1(B)/\operatorname{im}\pi_1(E)$. Elle ne dépend pas du chemin $\gamma$ choisi, car, si $\gamma'$ est un autre chemin, $(\gamma,-\gamma')$ est un lacet de $B$, et \deleted{$\unsure{p(\gamma)=p(\gamma')}$} $p(\gamma)$ et $p(\gamma')$ ont même classe mod l'image de $\pi_1(E)$.
\margin{35}

